\documentclass[11pt]{article}
\usepackage[margin=1in]{geometry}
\usepackage{hyperref}
\usepackage{enumitem}
\title{Environmental DNA for Biodiversity Monitoring: Promises, Pitfalls, and Statistical Frontiers\\
\large A Literature Review Generated with AI Assistance}
\author{LitReview AI \and A. L. Robles Fernández\\[4pt]
\small LitReview \textbar{} ai.litreview.org\\
\small Identifier: 2608.00002}
\date{2026-08-25}

\begin{document}
\maketitle

\begin{abstract}
Environmental DNA (eDNA) metabarcoding has transformed biodiversity monitoring by enabling species detection from water, soil, and air samples. This review synthesizes the state of the field across methodological and statistical dimensions. We examine laboratory protocols and contamination controls, bioinformatic pipelines for sequence assignment, and occupancy-style hierarchical models that translate detection counts into abundance estimates. We highlight persistent challenges: imperfect and heterogeneous detection, reference database gaps, quantitative interpretation of read counts, and spatial autocorrelation in study designs. Emerging directions include multi-trophic monitoring, eDNA time series for population trend estimation, and integration with autonomous sampling platforms.

\noindent\textbf{Keywords:} environmental DNA, metabarcoding, biodiversity monitoring, occupancy models.
\end{abstract}

\section*{Placeholder Notice}
This document is a placeholder entry in the LitReview repository. The full
review text will be generated with AI assistance from primary literature and
released after curation. This page establishes the identifier, metadata, and
distribution formats (LaTeX, PDF, HTML) for the entry.

\section*{Scope}
The forthcoming review will synthesize recent literature on the topic above,
covering its principal methods, findings, open challenges, and directions for
future work within the field of Ecology \& Evolution.

\section*{Data and Code Availability}
No new data or code are associated with this placeholder.

\bibliographystyle{plain}
\begin{thebibliography}{9}
\bibitem{litreview2026} LitReview Consortium (2026). \emph{LitReview: an arXiv-style repository of AI-generated literature reviews for the basic sciences.} ai.litreview.org.
\end{thebibliography}

\end{document}
